\begin{annexitem}{annex:p03.vaught-two.overview}{thm:p03-vaught-never-two}{The common never-two statement}
\begin{proof}
For a satisfiable complete theory the full construction, including finite
models and countable function languages, is
\cref{thm:p03-vaught-never-two-original}. An inconsistent theory has no
models, so its countable-model count is zero, also different from two.
Thus the overview needs no extra consistency exception.
\end{proof}
\end{annexitem}

\begin{annexitem}{annex:p03.vaught-two.types}{thm:p03-vaught-few-models-types}{Few models, countable types and isolated points}
\begin{proof}
Every complete finite-tuple type is realized in some model by compactness,
and then in a countable elementary submodel containing the realizing tuple
by downward L\"owenheim--Skolem. Choose representatives \(M_i\),
\(i\in J\subseteq\omega\), for the countable isomorphism types.
Consequently
\[
 S_m(T)=\bigcup_{i\in J}\{\operatorname{tp}_{M_i}(\bar a):
                         \bar a\in M_i^m\}.
\]
Each term is countable, as is the union; for \(m=0\) completeness gives
one type. The Stone topology has basic clopens
\([\varphi]=\{p:\varphi\in p\}\) and is compact by compactness of the
theory. Different types are separated by a formula.

Suppose a nonempty basic clopen contains no isolated type. Every nonempty
basic clopen within it contains two different types and therefore splits
into two nonempty basic clopens by some formula and its negation. Recursively
split along the finite binary strings. For each infinite binary sequence,
compactness gives a type in the nested clopens along its branch. Distinct
branches lie on opposite sides of a split and give different types. This
would inject \(2^\omega\) into the allegedly countable space. Hence
isolated types are dense. This is the atomistic Boolean-algebra step in
Vaught's notation.

If every type is isolated, the cover by their open singletons has a finite
subcover, so the space is finite. Conversely, in a finite space, for a fixed
type \(p\) choose one formula distinguishing it from each other type and
conjoin them. The resulting formula isolates \(p\). For the one-point
space use a tautology. This proves every assertion.
\end{proof}
\end{annexitem}

\begin{annexitem}{annex:p03.vaught-two.named}{thm:p03-vaught-named-finiteness}{Categoricity of a named-type theory and the formula-algebra injection}
\begin{proof}
In any complete countable-language theory \(Q\) with exactly one
denumerable model, every complete finite-tuple type must be isolated.
Otherwise the simultaneous omission construction of
Theorem~\ref{thm:p03-vaught-simultaneous-construction} supplies a
denumerable model omitting such a type, whereas compactness and downward
L\"owenheim--Skolem supply one realizing it. These models are not
isomorphic. By Lemma~\ref{thm:p03-vaught-few-models-types}, every
\(S_k(Q)\) is thus finite. A formula is determined modulo \(Q\) by the
types containing it: if formulas disagree, a realizing tuple witnesses a
type on which they disagree. There are at most \(2^{\abs{S_k(Q)}}\)
such subsets, so \(F_k(Q)\) is finite. Apply this to \(Q=T_p\).

The pointed theory is consistent by realizing \(p\). Every sentence in its
language is a formula of the original language with the named tuple
substituted, so completeness of \(p\) decides that sentence. It includes
\(T\). Suppose original formulas \(\varphi,\psi\) are inequivalent
modulo \(T\). The sentence
\(\exists\bar x\,\neg(\varphi\leftrightarrow\psi)\) is consistent
with \(T\), and completeness makes it a consequence of \(T\).
It holds in every model of \(T_p\). The formulas remain inequivalent
modulo \(T_p\), proving the injection. A nonisolated member of
\(S_m(T)\) forces that space, and hence \(F_m(T)\), to be infinite.
This contradicts the injection into finite \(F_m(T_p)\).
\end{proof}
\end{annexitem}

\begin{annexitem}{annex:p03.vaught-two.never}{thm:p03-vaught-never-two-original}{The complete original never-two argument}
\begin{proof}
First dispose of finite models. If a complete theory has a model \(A\)
of size \(d\), its sentence asserting exactly \(d\) elements holds
in every model. Suppose another model \(B\) is not isomorphic to \(A\).
There are finitely many bijections from \(A\) to \(B\). For each choose
an atomic instance that it fails to preserve or reflect. The finite
conjunction of these instances, together with the full finite diagram on
the symbols involved, gives an existential sentence true in \(A\) but
false in \(B\). This contradicts completeness. Hence all models are
isomorphic and the spectrum is one.

Now suppose all models are infinite and there are exactly two countable
isomorphism types. Every \(S_k(T)\) is countable by
Lemma~\ref{thm:p03-vaught-few-models-types}. The identity-interpretation
case of Theorem~\ref{thm:p03-vaught-saturated-existence} supplies a
denumerable saturated model \(S\).

Some \(S_m(T)\) has a nonisolated member \(p\). Indeed, if all types
were isolated, every denumerable model would be atomic and the original
atomic uniqueness theorem
(Theorem~\ref{thm:p03-vaught-atomic-uniqueness}) would make any two
isomorphic. This contradicts the hypothesis of two types.
Use Theorem~\ref{thm:p03-vaught-simultaneous-construction} to obtain
a denumerable model \(O\) omitting \(p\). It is not saturated, because
\(p\) is a complete finite-tuple type over the empty parameter set.
Thus \(O\not\cong S\), and these exhaust the two types. Every
denumerable model realizing \(p\) must be isomorphic to \(S\), and
therefore saturated.

Let \((M,\bar a)\) and \((N,\bar b)\) be arbitrary denumerable
models of \(T_p\). Their reducts realize \(p\), so both are saturated.
The map \(\bar a\mapsto\bar b\) is a well-defined finite elementary
map: equalities and every other formula agree because both tuples realize
\(p\). Enumerate both domains. Extend the map alternately to the next
unused element of \(M\) and of \(N\), using saturation over the finite
tuple already mapped. At a forth step, transport the complete type of the
chosen element over that tuple to the image tuple. Every finite part is
consistent there by elementarity of the finite map; saturation realizes the
whole type. Its inequalities keep the new image distinct from previously
mapped elements. The back step is the same argument in the reverse
direction. The union is bijective and elementary because every formula
uses finitely many arguments and is decided at some finite step.
It sends \(\bar a\) to \(\bar b\), so is an isomorphism of the
pointed structures. This is Vaught's Lemma~4.5 argument with the initial
tuple prescribed, and proves \(T_p\) denumerably categorical.
Lemma~\ref{thm:p03-vaught-named-finiteness} now contradicts the
nonisolation of \(p\). Exactly two types are impossible.

For a countable language with functions, replace each function by its graph
relation and add totality and uniqueness axioms. Every original structure
has a unique graph expansion, and every model of the graph axioms has a
unique function reconstruction. Isomorphisms correspond in both directions.
Translate terms by finitely many existentially quantified graph witnesses;
totality and uniqueness prove, by formula induction, that truth is preserved.
The translations of a complete theory together with the graph axioms are
complete, since every relational sentence translates back to a function
sentence on those models. The language stays countable and model cardinalities
stay unchanged. Its supposed two countable models would contradict the
relational result already proved.
\end{proof}
\end{annexitem}

\begin{annexitem}{annex:p03.vaught-two.examples}{thm:p03-vaught-finite-counts}{Completeness, all cut cases and the finite-language construction}
\begin{proof}
We first record the elementary order construction used throughout. Two
countable dense orders without endpoints, with the same finite list of dense
colors, are isomorphic by back-and-forth. A finite color-preserving partial
order isomorphism can be extended to any next element: its required image
lies in an open interval, an initial ray, or a final ray, and the needed
color occurs there. Choose a fresh image; density supplies infinitely many.
Alternate the two enumerations. The same finite-extension argument, without
surjectivity, proves preservation of formulas by induction on quantifiers:
extend a finite map to an existential witness and use induction on its
matrix; universal quantifiers follow by negation.

In \(T_{r+2}\), a sentence names only finitely many \(c_i\). Take an
initial segment containing them all. Between its successive constants,
below its first constant and above its last, the underlying colored order
has exactly the interval structure just described. Sending the finite list
of constants to their counterparts starts the witness-extension argument.
Hence all models agree on the sentence, proving completeness of the entire
countable-language theory.

For its countable models put \(B=\{x:(\forall i)c_i<x\}\), where this
is a metalinguistic definition. If \(B\) is empty, the constants are
cofinal. If it is nonempty without a least member, \(B\) itself is a
dense colored order without endpoints. If it has a least member \(b\),
the color of \(b\) is one of \(r\) choices; \(B\setminus\{b\}\)
is again such an order. Below \(B\), split into
\(( -\infty,c_0)\), the constant points, and the intervals
\((c_i,c_{i+1})\). Each open piece has the same countable colored order
type. Isomorphisms of these pieces, together with the prescribed constants
and, where present, \(b\), assemble to an isomorphism of any two models
in the same case. Different cases, or different colors of \(b\), cannot
be isomorphic. All cases exist: concatenate the indicated countable dense
colored pieces, adjoining the specified constant points and optional least
upper-bound point. Every color remains dense and there are no endpoints.
This gives exactly \(r+2\) types. For one model, use the uncolored dense
order without constants and the first back-and-forth argument. The negative
two case is Theorem~\ref{thm:p03-vaught-never-two-original}.

For \(T'_{r+2}\), its quotient is such a colored order with distinguished
points \(C_i\). In a countable model all nondistinguished fibers are
countably infinite, while the fiber over \(C_i\) has size \(i\).
Thus the quotient classification just proved lifts to the full structure:
combine its quotient isomorphism with arbitrary fiber bijections. Conversely,
an isomorphism preserves \(R\), every fiber size, every color, and the
quotient order. It must fix the distinguished quotient points by their
unique sizes. Exactly the same \(r+2\) cases occur, and replacing each
quotient point by its required finite or countably infinite fiber constructs
them all.

It remains to prove completeness of this finite-language theory rather than
infer it from its model count. Fix a positive finite game length \(q\).
Mark in both quotients the classes \(C_1,\ldots,C_q\); match them by
their unique sizes. Maintain a partial correspondence of the classes visited
by the chosen elements, preserving their order and color and their position
relative to the marked classes. Inside a visited fiber maintain an injective
partial map. A move in a marked fiber has a free matching element whenever
the source move is new, because those fibers have equal size. A move in an
unmarked visited fiber has a free matching element because both fibers have
at least \(q+1\) members and at most \(q\) total moves occur.
For a new unmarked class, choose a class of the same color in the prescribed
open interval or ray determined by the marked and visited classes.
Density of that color and finiteness of the excluded classes supply a fresh
one, necessarily unmarked; it has at least \(q+1\) members. This strategy
works for moves in either structure and preserves all atomic relations.
Induction on formulas of quantifier rank at most the remaining game length
proves agreement on those formulas: for an existential quantifier, play its
witness and apply the induction to the matrix; Boolean connectives preserve
agreement. Every sentence has finite quantifier rank, so every two models
of \(T'_{r+2}\) are elementarily equivalent. This proves completeness.
The notation \(C_i\) in this game is external marking, not an illicit
addition to the finite language.
\end{proof}
\end{annexitem}

\begin{annexitem}{annex:p03.vaught-two.equivalence}{thm:p03-vaught-equivalence-example}{Robinson's equivalence-relation example}
\begin{proof}
A finite partial bijection respecting \(R\) extends to a new element
in the corresponding equivalence class because that class is infinite.
If it is the first visited member of its class, choose the corresponding
as yet unvisited class. There are exactly two classes on each side, so this
is possible. The symmetric extension proves equality of all first-order
sentences by induction on quantifiers. The theory is complete.

In a denumerable model both classes are countably infinite and arbitrary
class bijections give an isomorphism. In a model of size \(\aleph_1\)
at least one class has that size, because a union of two countable sets
is countable. The other infinite class has size either \(\aleph_0\)
or \(\aleph_1\); there is no cardinal strictly between them. Both
possibilities exist and the unordered pair of sizes completely determines
the isomorphism type. This proves the stated two uncountable models.
\end{proof}
\end{annexitem}

\begin{annexitem}{annex:p03.vaught-two.variants}{thm:p03-vaught-cardinal-boundaries}{Counting bounds, perfect sets, CH and a two-model infinitary sentence}
\begin{proof}
For structures on \(\omega\), each finite-arity relation is coded by a
subset of a countable set, each function by a member of a countable product
of \(\omega\), and each constant by a natural number. A countable
language therefore has at most \(2^{\aleph_0}\) such codes. Every
denumerable structure can be transported to \(\omega\). Finite domains
add at most continuum many codes, using their analogous countable-coordinate
spaces. Passing to isomorphism classes cannot increase this cardinal bound.

Each coding component is a countable product of discrete Polish spaces.
The countable disjoint union over finite domain sizes and the domain
\(\omega\) is also Polish: give distinct components distance one and use
a bounded complete compatible metric within each component. It has a
countable base. A nonempty perfect closed
subset contains a copy of \(2^\omega\): recursively choose two disjoint
small neighborhoods meeting it, each with closure contained in its parent,
and with diameters tending to zero in a complete compatible metric.
There is room for two choices since it has no isolated point. Completeness
and closedness give one limit for each branch; disjoint siblings make the
limits distinct. The preceding coding bound gives the reverse cardinal
inequality. If the members are pairwise nonisomorphic, they supply continuum
many distinct isomorphism types. Under CH, every uncountable cardinal at
most continuum equals \(\aleph_1=2^{\aleph_0}\), proving the mere
cardinal alternative. No part of that cardinal calculation selects a closed
subset or constructs the neighborhood tree, so it supplies no perfect-set
witness.

Let \(\tau\) be the countable conjunction of the first-order axioms of
\(T_3\), and put
\[
 \sigma=\tau\wedge\exists x\,\bigwedge_{i<\omega}(c_i<x).
\]
This is a legitimate \(L_{\omega_1\omega}\)-sentence: the displayed
conjunction has only the one free variable \(x\). It excludes exactly
the cofinal model. By Theorem~\ref{thm:p03-vaught-finite-counts}, the
remaining cases, upper cut without a least element and upper cut with a
least element, each have one denumerable isomorphism type. Both reducts
satisfy the same complete first-order theory \(T_3\). Thus \(\sigma\)
has exactly two models in that cardinal without contradicting never-two:
its countable model class is an infinitary restriction of the first-order
class, not the full model class of a complete first-order theory.
\end{proof}
\end{annexitem}
